\section{Conclusion}

This paper introduced \toolname, a fuzzer designed to detect
under-invalidation defects in modern browser layout engines.
Unlike fuzzers that compare different browsers or browser versions,
\toolname uses self-differential testing to compare a browser's incremental
layout with its own from-scratch layout. This comparison provides a practical
oracle for detecting stale layout state after style changes without assuming
that another implementation provides ground truth.

\toolname combines this oracle with grammar-directed generation, large DOM and
CSS test cases, automated minimization, and structural clustering. These
components enable \toolname to detect layout defects while reducing the
complexity and volume of the test cases that developers must review. The
evaluation showed that \toolname reproduced known under-invalidation bugs,
tested hundreds of candidates per minute in Chrome and Firefox, and
substantially reduced review volume through minimization and clustering.
Furthermore, \toolname identified \NumBugsNovel{} novel bugs that browser
developers accepted, of which \NumBugsFixed{} have been fixed.

These findings establish under-invalidation as a productive target for browser
fuzzing and demonstrate that self-differential testing can detect these
defects effectively. As browser layout engines incorporate increasingly
complex features, \toolname can help developers identify subtle correctness
defects earlier and improve web-platform reliability. The same approach may
also apply to systems that rely on invalidation, including graphical user
interface frameworks, spreadsheet-calculation pipelines, and incremental
build tools.
